\chapter{PONDER-05: DC-Biased Quartz Thruster}
\label{ch:ponder_05}

The PONDER-05 configuration accesses Axiom 4's nonlinear dielectric regime by applying a static $30\text{ kV}$ DC bias across a quartz cylinder, overlaid with a modest $500\text{ V}$ AC perturbation at $50\text{ kHz}$. The DC bias holds the material at $68.7\%$ of the kinetic yield voltage ($V_{yield} = 43.65\text{ kV}$), where the effective vacuum permittivity has already collapsed by $27.4\%$ from its zero-field value while the effective capacitance has risen by $37.7\%$.

\section{The Nonlinear Operating Regime}

The universal saturation kernel $S$ (Axiom~4) governs \emph{two} measurable quantities with opposite behaviours:
\begin{equation}
    \varepsilon_{eff}(V) = \varepsilon_0 \cdot \underbrace{\sqrt{1 - \left(\frac{V}{V_{yield}}\right)^{\!2}}}_{S \;\to\; 0}
    \qquad\qquad
    C_{eff}(V) = \frac{C_0}{\underbrace{\sqrt{1 - \left(\frac{V}{V_{yield}}\right)^{\!2}}}_{S \;\to\; 0}} \;\to\; \infty
\end{equation}
where $V_{yield} = \sqrt{\alpha} \cdot V_{snap} \approx 43.65\text{ kV}$. At $V_{DC} = 30\text{ kV}$:
\begin{equation}
    S\big|_{30\text{kV}} = \sqrt{1 - \left(\frac{30}{43.65}\right)^{\!2}} = \sqrt{1 - 0.4726} \approx 0.726
\end{equation}

The constitutive permittivity \emph{drops} to $72.6\%$ of its zero-field value (the vacuum becomes a weaker dielectric), while the measured capacitance \emph{rises} to $1/0.726 \approx 137.7\%$ (the node absorbs more displacement current under strain).

\section{DC Cross-Term Amplification}

When a small AC signal $V_{ac}\sin(\omega t)$ rides on a large DC bias $V_{DC}$, the total squared field contains a cross term:
\begin{equation}
    |E|^2 = (V_{DC} + V_{ac}\sin\omega t)^2 = V_{DC}^2 + 2V_{DC}V_{ac}\sin\omega t + V_{ac}^2\sin^2\omega t
\end{equation}

The $2V_{DC}V_{ac}$ cross term provides \textbf{linear amplification} of the AC signal. For $V_{DC} = 30\text{ kV}$ and $V_{ac} = 500\text{ V}$:
\begin{equation}
    \text{Amplification} = \frac{2V_{DC}}{V_{ac}} = \frac{2 \times 30{,}000}{500} = 120\times
\end{equation}

This $120\times$ cross-term amplification transforms a barely detectable AC-only thrust into a $469\,\mu\text{N}$ signal---nearly $500\times$ above the $1\,\mu\text{N}$ torsion balance detection floor.

\section{Rarefaction Wave Inversion}

Near $V_{yield}$, the nonlinear $\varepsilon_{eff}(V)$ curve creates a velocity asymmetry. The local phase velocity of acoustic perturbations depends on the instantaneous voltage:
\begin{equation}
    c_{phase}(V) \propto \frac{1}{\sqrt{\varepsilon_{eff}(V)}} \propto \left(1 - \left(\frac{V}{V_{yield}}\right)^2\right)^{1/4}
\end{equation}

At the positive peak of the AC cycle ($V_{DC} + V_{ac}$), the phase velocity is \textit{slower} than at the negative peak ($V_{DC} - V_{ac}$). This velocity asymmetry causes the wave crest to lag relative to the trough---the exact opposite mechanism of standard forward acoustic shockwaves.

Rather than steepening into a compressive shock front, the topological wave stretches out into a \textbf{rarefaction fan} (a backward-breaking wave). Because the steepening evaluates to a fraction less than 1, any nonlinear excess forward thrust previously derived from this mechanism is inverted. Further research is required to determine if a continuous rarefaction wave drawn through a dielectric lattice produces unidirectional macro-thrust, or if it acts as a reverse-thrust drag brake on the PONDER system.

\section{The Quartz Cylinder}

\begin{center}
\begin{tabular}{|l|c|c|}
\hline
Parameter & Value & Rationale \\ \hline
Material & X-cut $\alpha$-quartz & Piezoelectric, ultralow loss \\
Diameter & 50 mm & Standard resonator size \\
Length & 50 mm & $\lambda/2$ at 50 kHz \\
$\varepsilon_r$ & 4.5 & Low permittivity $\to$ high field penetration \\
$\tan\delta$ & $10^{-5}$ & Negligible thermal dissipation \\
Electrodes & Sputtered Cr/Au & Ohmic contact, no air gap \\
AC frequency & 50 kHz & Piezoelectric resonance \\
AC amplitude & 500 V RMS & Standard ultrasonic driver \\
DC bias & 30 kV & $68.7\%$ of $V_{yield}$ \\
\hline
\end{tabular}
\end{center}

The quartz is driven at its fundamental thickness-mode resonance ($50\text{ kHz}$ for a $50\text{ mm}$ cylinder). The AC drive uses a standard ultrasonic piezo driver. The $30\text{ kV}$ DC bias is supplied by a laboratory HV power supply through a large blocking inductor ($> 1\text{ H}$) to prevent AC current from flowing into the DC supply.

\section{The Mineral Oil Dielectric Bath}

The quartz cylinder operates submerged in degassed mineral oil, providing three critical engineering advantages simultaneously.

\subsection{Corona Suppression}

At $30\text{ kV}$ across a $50\text{ mm}$ gap, the applied field is $E_{app} = 0.60\text{ MV/m}$. Mineral oil has a breakdown field of $12\text{ MV/m}$, providing a $20\times$ safety margin against arcing. In contrast, air at STP has only a $5\times$ margin, and surface tracking along the quartz cylinder can produce ion-wind artifacts indistinguishable from thrust.

\subsection{Impedance Matching}

The mineral oil layer ($\varepsilon_r = 2.2$, $Z \approx 254\,\Omega$) acts as a quarter-wave impedance transformer between quartz ($Z \approx 178\,\Omega$) and free space ($Z_0 = 377\,\Omega$):
\begin{equation}
    \text{Reflected power:} \quad
    \begin{cases}
        \text{Quartz} \to \text{Vacuum (direct):} & 12.9\% \\
        \text{Quartz} \to \text{Oil} \to \text{Vacuum:} & \sim 3.4\%
    \end{cases}
\end{equation}

This $3.7\times$ reduction in reflected acoustic power ensures that the generated thrust wave efficiently radiates into the surrounding medium.

\subsection{Thermal Management}

Quartz dissipates $< 0.001\text{ mW}$ at these operating conditions ($\tan\delta = 10^{-5}$). The oil bath provides additional convective cooling ($h \approx 100\text{ W/m}^2\text{K}$), rendering thermal artifacts entirely negligible. The temperature rise of the quartz surface is $< 0.001$\textdegree C even after hours of continuous operation.

\section{Predicted Thrust Profile (Linear Base)}

Because the acoustic shockwave mechanism mathematically inverts into a rarefaction wave under the correct structural denominator operator, the purely nonlinear excess thrust cannot be predicted by standard compressive fluid dynamics. The table below represents only the linear cross-term amplifications as a function of the local permittivity modulation ($\varepsilon_{eff}/\varepsilon_0$).

\begin{center}
\begin{tabular}{|c|c|c|}
\hline
$V_{DC}$ (kV) & Cross-Term $F_{lin}$ ($\mu$N) & $\varepsilon_{eff}/\varepsilon_0$ \\ \hline
0   & 0    & 1.000 \\
10  & 156  & 1.027 \\
20  & 313  & 1.125 \\
30  & 469  & 1.377 \\
40  & 626  & 2.493 \\
43  & 672  & 5.923 \\ \hline
\end{tabular}
\end{center}

The linear cross-term currently serves as the only mathematically sound bound for macroscopic thrust generation, pending a rigorous re-derivation of the macroscopic rarefaction wave integral using the Maxwell Stress Tensor.

\begin{figure}[h]
    \centering
    \includegraphics[width=1.0\textwidth]{ponder05_saturation_curves.png}
    \caption{Left: the dual behaviour of the Axiom~4 saturation kernel --- $\varepsilon_{eff}$ collapses toward zero while $C_{eff}$ diverges as voltage approaches $V_{yield} = 43.65$~kV. The 30~kV DC bias point is marked. Right: the linear cross-term thrust scales linearly with DC bias, providing a $469\,\mu$N signal at 30~kV --- $469\times$ above the torsion balance detection floor.}
    \label{fig:ponder05_saturation}
\end{figure}

\section{Bill of Materials}

\begin{center}
\begin{tabular}{|l|c|c|}
\hline
Item & Quantity & Est.\ Cost \\ \hline
Custom X-cut quartz cylinder (50 mm) & 1 & \$400 \\
HV DC power supply (0--50 kV, 1 mA) & 1 & \$1{,}200 \\
Ultrasonic piezo driver (50 kHz, 500 V) & 1 & \$300 \\
HV blocking inductor ($> 1$ H) & 1 & \$150 \\
Mineral oil (transformer grade, 5 L) & 1 & \$50 \\
PTFE test chamber / oil bath & 1 & \$200 \\
HV cabling, connectors, feedthroughs & lot & \$300 \\
Torsion balance (see Chapter~5) & 1 & \$400 \\ \hline
\textbf{Total} & & \textbf{$\sim$\$3{,}000} \\ \hline
\end{tabular}
\end{center}
